\documentclass[11pt]{article}

\usepackage[margin=1in]{geometry}
\usepackage{amsmath,amssymb,amsthm}
\usepackage{booktabs}
\usepackage{xcolor}
\usepackage{listings}
\usepackage[colorlinks=true,linkcolor=blue,urlcolor=blue]{hyperref}

\lstset{
  basicstyle=\ttfamily\small,
  breaklines=true,
  frame=single,
  backgroundcolor=\color{black!4},
  columns=fullflexible,
  keepspaces=true,
}

\newcommand{\Real}{\operatorname{Re}}
\newcommand{\Imag}{\operatorname{Im}}
\newcommand{\atil}{\tilde{a}}
\newcommand{\code}[1]{\texttt{#1}}

\title{\textbf{The \texttt{liftcomplex} obstruction:\\
why auxiliary--variable lifting fails for the bubble at imaginary $\mu$}}
\author{Tropical Monte Carlo (FINAL4) --- planAU Phase-0 finding}
\date{2026-06-29}

\begin{document}
\maketitle

\begin{abstract}
The auxiliary--variable \emph{lift} (v3 engine, \code{ProcessSectorLifted}) folds a
large kinematic coefficient into the Newton polytope so the tropical sampling map
matches a hierarchical integrand. It is the intended fix for the noisy
imaginary--$\mu$ ``principal series'' (FLAG~B). On the actual bubble integrand at
imaginary $\mu$ the lift refuses with \code{TropicalEval::liftcomplex} on every
cone. This note explains \emph{exactly why}: the lift introduces a real
\emph{domain inequality} that the ordinary (unlifted) algorithm does not have, and
that inequality is governed by an effective exponent $\atil$ built from
\emph{both} the polynomial exponents $B_j$ \emph{and} the monomial exponents
$m_i$. The \code{SplitRealImag} complex--exponent mode removes the imaginary part of
$B_j$ before the inequality is formed, but \emph{not} that of $m_i$; the bubble at
imaginary $\mu$ is the first integrand with complex \emph{monomial} exponents, so
$\atil$ stays complex and the lift aborts. The flattening (change of variables) is
\emph{not} the broken step. A one--character experiment confirms the diagnosis and
points to a localized fix.
\end{abstract}

\section{Setup and notation}

Each presector integrand has the form
\begin{equation}
  I \;=\; \int_{[0,\infty)^n}\; \prod_{i=1}^{n} x_i^{\,m_i}\;
          \prod_{j} P_j(x)^{\,B_j}\; \mathrm{d}^n x ,
  \label{eq:integrand}
\end{equation}
with \emph{monomial exponents} $m_i$ (on the bare powers $x_i^{m_i}$) and
\emph{polynomial exponents} $B_j$ (on the polynomials $P_j$). For the bubble these
exponents are functions of the propagator dimensions $\delta_\pm = \tfrac{3}{2}\pm\mu$
and the regulator $d_\epsilon$. For \emph{real} $\mu$ they are real; for
\emph{imaginary} $\mu$ (the principal series) they are \textbf{complex}, and --- this
is the crux --- \emph{both} families $m_i$ and $B_j$ acquire nonzero imaginary parts.

\medskip
\noindent\textbf{Concrete instance.} For \code{bubblepm} presector~1 at $\mu = 3i/2$,
\[
  \Imag(m) \;=\; \bigl\{\,-\tfrac32,\ \tfrac32,\ -\tfrac32,\ \tfrac32,\ 3\,\bigr\},
  \qquad
  \Imag(B) \neq 0 \ \text{as well.}
\]

\section{The ordinary algorithm never \emph{decides} on an exponent value}

The tropical decomposition triangulates the dual fan into simplicial cones. In each
cone it applies a \emph{monomial change of variables}
\begin{equation}
  x_i \;=\; \prod_{a} y_a^{\,(M)_{ia}},
  \label{eq:cov}
\end{equation}
where the integer matrix $M$ is fixed by the cone's rays --- i.e.\ by the
\emph{support} (the set of exponent vectors), \emph{not} by the exponent
\emph{values}. Substituting \eqref{eq:cov} into \eqref{eq:integrand} and clearing the
leading monomial from each polynomial is the \emph{flattening} step
(\code{ProcessSector}). The exponents simply ride along as numbers; the integrand is
evaluated as $\exp\!\big(B_j \log P_j\big)$, which is well defined for \emph{any}
complex $B_j$.

\begin{quote}\itshape
Nothing in the unlifted pipeline asks a yes/no question about an exponent. Complex
exponents flow through untouched --- which is precisely why the unlifted scan of the
imaginary--$\mu$ masses completed normally.
\end{quote}

\section{The lift adds a real \emph{domain inequality}}

To tame a hierarchical coefficient $C\,x^{\alpha}$ the lift rewrites it as
$c\,z^{k} x^{\alpha}$, promotes the anchor $z$ to a genuine integration variable, and
pins it with a delta function:
\begin{equation}
  C\,x^{\alpha}\ \longrightarrow\ c\,z^{k}x^{\alpha},
  \qquad
  I \;=\; \int_{[0,\infty)^{\,n+1}} (\cdots)\;\delta(z-z_0)\,\mathrm{d}z\,\mathrm{d}^n x ,
  \qquad z_0 = |C|^{1/k}.
\end{equation}
One flattens this $(n{+}1)$--dimensional problem and then \textbf{resolves the
delta}: a pivot variable $y_p$ is chosen and eliminated, folding the integral back to
$n$ dimensions. That fold--back is the new ingredient. Solving $\delta(z-z_0)$ for the
pivot turns the sector's domain into a \emph{cut region}, emitted as an indicator
\begin{equation}
  \mathrm{Boole}\!\left[\, \log y_p^{\star} \le 0 \,\right],
  \qquad
  \log y_p^{\star}
   \;=\; \frac{\log z_0 - \sum_a \mathrm{ic}_a \,\log y_a}{m_p},
  \label{eq:indicator}
\end{equation}
together with a classification of the sector (\code{HasConstantTerm},
\code{EmptyDomain}, divergent/convergent). The coefficients in \eqref{eq:indicator}
and the classification are read off the \textbf{effective exponents} $\atil$ left
after the delta is resolved.

\paragraph{Why $\atil$ must be real.} Equation~\eqref{eq:indicator} is an
\emph{inequality}: it carves $[0,1]^d$ into a region that contributes and a region
that does not. ``Is $(2-3i)\le 0$?'' is meaningless. Hence the lift requires
$\atil\in\mathbb{R}$. The engine guards this explicitly; when every admissible pivot
yields a complex $\atil$ it aborts:
\begin{lstlisting}
TropicalEval::liftcomplex = "ProcessSectorLifted: cone `1` -- all candidate
  pivots produce complex atilde; cannot emit a real-valued domain indicator."
\end{lstlisting}

\paragraph{Crucially, $\atil$ depends on \emph{both} exponent families.} Under the
change of variables \eqref{eq:cov} the bare factor $\prod_i x_i^{m_i}$ becomes
\begin{equation}
  \prod_i x_i^{m_i}
  \;=\;\prod_a y_a^{\,\sum_i m_i (M)_{ia}},
  \label{eq:monoshift}
\end{equation}
so the new variables $y_a$ inherit exponents $\sum_i m_i (M)_{ia}$ that are
\textbf{linear in the monomial exponents $m_i$}. These combine with the
polynomial--derived exponents; the delta resolution then produces
\begin{equation}
  \atil \;=\; \atil\big(m,\,B,\,\text{cone geometry}\big),
  \qquad\text{linear in } m \text{ and } B .
  \label{eq:atilde}
\end{equation}
If \emph{any} contributing exponent is complex, $\atil$ is complex and the
inequality~\eqref{eq:indicator} is ill posed.

\section{\texttt{SplitRealImag}: the intended fix, wired to only half the inputs}

The engine's mechanism for complex exponents is \code{ComplexExponentMode ->
"SplitRealImag"}: build the (real) domain map from the \emph{real parts} of the
exponents, and peel the \emph{imaginary parts} off into a pure oscillatory phase
\begin{equation}
  \exp\!\Big(\, i\,\textstyle\sum_j \Imag(B_j)\,\big(\log Q_j + \mathrm{MonoFactorLog}_j\big)\Big)
\end{equation}
that multiplies the integrand. A phase does not move the boundary, so
\eqref{eq:indicator} stays real. This is correct and is validated for complex
\emph{polynomial} exponents (\code{cc\_21b}, \code{cc\_45}).

The defect is that the split is applied to only one of the two inputs to
\eqref{eq:atilde}. In \code{EvaluateTropicalMC} (file \code{tropical\_eval.wl},
$\approx$ lines 3772 and 3814):
\begin{lstlisting}
imBList  = Im[integrandSpec["PolynomialExponents"]];      (* polynomial only *)
hasImagB = AnyTrue[imBList, (! TrueQ[PossibleZeroQ[#]]) &];
...
ProcessSectorLifted[
  MapAt[Re, liftedSpec, {Key["PolynomialExponents"]}],     (* Re of B only;  *)
  ...                                                       (* m left complex *)
  "Eps" -> eps]
\end{lstlisting}
Only $\Real(B_j)$ is taken; the monomial exponents $m_i$ are passed through
\textbf{complex}. By \eqref{eq:monoshift}--\eqref{eq:atilde}, $\atil$ is therefore
\emph{still} complex --- now from the monomial half --- and the \code{liftcomplex}
guard fires on every cone.

\begin{quote}\itshape
The bubble at imaginary $\mu$ is the first integrand anyone has lifted whose
\emph{monomial} exponents are complex; all prior complex tests had complex $B_j$ with
real $m_i$, so the incomplete split was invisible upstream.
\end{quote}

\section{Empirical confirmation}

On the \code{bubblepm} $\mu=3i/2$ lifted spec (presector~1; $n{=}5$, so $611$
six--dimensional cones), processing each cone two ways:

\begin{center}
\begin{tabular}{lcc}
\toprule
What is fed to \code{ProcessSectorLifted} & complex $\atil$? & Result \\
\midrule
$\Real$ of \emph{polynomial} exponents only \emph{(what the engine does)} & yes & \code{liftcomplex} --- \textbf{every cone fails} \\
$\Real$ of \emph{both} monomial \emph{and} polynomial exponents & no & \textbf{all 611 cones succeed} \\
\bottomrule
\end{tabular}
\end{center}

\noindent So the obstruction is \emph{solely} the un--split monomial exponents, and
making $\atil$ real is sufficient to let the decomposition through. (For reference,
the unlifted path and an independent \code{NIntegrate} of \eqref{eq:integrand} both
evaluate this point without issue, giving $\approx 1.27\times10^{-14} +
1.40\times10^{-15}\,i$.)

\section{What is and is not broken}

\begin{itemize}
  \item \textbf{Flattening / change of variables --- NOT broken.} It is geometric
  (fixed by the cone rays), shared verbatim with the unlifted path, and handles
  complex exponents fine. The imaginary parts \emph{originate} in the flattened
  exponents via \eqref{eq:monoshift}, but storing complex numbers is harmless.
  \item \textbf{Delta--resolution / pivot / domain indicator (lift--only) ---
  broken.} This step needs $\atil\in\mathbb{R}$ to form the
  inequality~\eqref{eq:indicator} and to classify the sector; with complex $m_i$ it
  refuses (\code{liftcomplex}).
\end{itemize}

\section{The fix (sketch)}

Extend \code{SplitRealImag} to treat the monomial exponents on the same footing as
the polynomial exponents:
\begin{enumerate}
  \item \textbf{Decomposition side.} Take $\Real(m_i)$ as well as $\Real(B_j)$ before
  \code{ProcessSectorLifted}, so $\atil$ in \eqref{eq:atilde} is real and the domain
  indicator is well posed. \emph{(Shown sufficient above: $611/611$ cones.)}
  \item \textbf{Phase side.} Carry $\Imag(m_i)$ into the \emph{same} oscillatory phase
  that already transports $\Imag(B_j)$. Using \eqref{eq:cov},
  \[
    \sum_i \Imag(m_i)\,\log x_i
    \;=\; \sum_a \Big(\textstyle\sum_i \Imag(m_i)\,(M)_{ia}\Big)\log y_a ,
  \]
  i.e.\ the monomial imaginary parts contribute
  $\sum_i \Imag(m_i)(M)_{ia}$ to the phase coefficient of $\log y_a$ in every sector
  (plus the matching anchor/$\mathrm{MonoFactorLog}$ constant from the lift). This is
  the part requiring care: it must be \emph{numerically correct}, not merely
  non--crashing, and must reduce to the current behaviour when $\Imag(m_i)=0$ (so
  real--$\mu$ and prior complex--$B$ outputs stay byte--identical).
\end{enumerate}
The change is local to the lifted \code{SplitRealImag} path and its codegen (both the
convergent route, e.g.\ \code{bubblepm}, and the divergent IBP / subtraction route,
e.g.\ \code{bubblepp}); flattening is untouched. Validation should mirror
\code{cc\_21b}: (a) lifted \code{SplitRealImag} reproduces \code{NIntegrate} of
\eqref{eq:integrand} to $<1$--$2\%$; (b) dropping the new monomial phase term gives a
\emph{wrong} answer (so it is load--bearing); (c) real--$\mu$ results are unchanged.

\end{document}
